\scanpage{104}\mspage{95}
% archive sheet number: 104
a) \unsure{préfigures|Préfigures} \uline{\unsure{fermées|fermée}} $0$-connexes \inserted{$=$ figures $0$-connexes} -- déterminées par les éléments de $\mathrm{Drap}^{\unsure{*|\times}}(\unsure{L|\mathcal L})$, cf plus haut

b) Figures \deleted{\unsure{f}} ouvertes $0$-connexes : \deleted{\unsure{cas de}} elles sont données par $\deleted{\unsure{D}}T\in\mathrm{Drap}^{*}(L)$ \uline{$\mathrm{card}\,\deleted{\unsure{D}}T\geq 2$}, et une partie de
\[
\begin{aligned}
 \partial F_T\overset{\text{def}}{=}\unsure{\partial|\delta}T
 &=\{\text{ensemble formé du plus}\\
 &\qquad\text{petit et du plus grand}\\
 &\qquad\text{élément de }T=F_0\}
\end{aligned}
\]
\[
 \Phi=F\setminus\partial\unsure{F|T}
\]
\margin{\uline{NB} On a finalement identifié les $P_a$ aux $a$, i.e. $\unsure{M|\mathcal M}_0$ à $L$ \ldots}

c) Cas général : on se donne une partie $\deleted{\unsure{\beta}}\alpha$ de $\partial F$, et on prend
\[
 \Phi=F\setminus\alpha
\]

\uline{Ex} Segments fermé, semi-ouvert, ouvert (\deleted{\unsure{cas}} $\mathrm{card}\,T=2$)

\deleted{\unsure{B|D}}Soit $\text{PréFig conn}(M)$ l'ensemble des préfigures connexes, qu'on vient d'expliciter, identifié aux couples $(T,\alpha)$, $T\in\mathrm{Drap}^{*}(L)$, $\alpha\in\partial T$ (\deleted{frontière}''bord de la figure \inserted{connexe} $F_T$ ; \uline{NB} si $F_T=\{P_x\}$, on \deleted{\unsure{aura}} \inserted{prendra} $\partial F_T=\varnothing$).
% slanted
\margin{On pose $F=\overline{\Phi}$ (plus petite figure \deleted{fermée} \unsure{contenant} la préfigure), $\partial F=\partial\Phi$}

\deleted{On va définir \unsure{sur|dans}}
\[
 \deleted{\text{Préfigconn}(M)}
\]
\deleted{une relation d'ordre} $\deleted{\inserted{<}}$, \deleted{en écrivant}
\[
 \deleted{\Phi<\Phi'\Longleftrightarrow}
 \left\{\begin{aligned}
 &\deleted{\text{a) pour tout }X\in\Phi,X'\in\Phi',\ \text{on a}}\\
 &\deleted{X\underset{\mathrm{pos}}{<}X'\ \text{\unsure{au sens}}}\quad\deleted{\text{cf plus bas}}\\
 &\deleted{\text{\unsure{b)}}}\quad\deleted{-\ \text{\unsure{cel|cela|ceci} implique }\mathrm{card}\,\partial\Phi\cap\partial\Phi'\leq1}\\
 &\deleted{\text{b) Si }x\in\partial\Phi\cap\partial\Phi',\ \text{on a }x\notin\Phi,x\notin\Phi'}
 \end{aligned}\right.
\]
\margin{\uline{NB} b) \unsure{indique|implique} que si $\Phi$ de dim $0$ i.e $\Phi=\{P_a\}$ \inserted{alors $\unsure{\partial|\delta}$}}
